\documentclass[10pt,a4paper]{amsart}
\usepackage[margin=19mm]{geometry}
\usepackage{amsmath,amssymb,mathtools}
\usepackage[scr=rsfs,cal=euler]{mathalfa}
\usepackage{xfrac}
\usepackage[unicode,hidelinks,bookmarks=false]{hyperref}
\newcommand{\ie}{i.e.\ }
\newcommand{\cf}{cf.\ }
\newcommand{\resp}{respectively\ }
\newcommand{\Subsection}{\subsection}
\providecommand{\nobreakdash}{\nobreak\hbox{-}}
\setlength{\emergencystretch}{2em}
\multlinegap0pt
\usepackage{amssymb}
\usepackage{enumerate}
\newtheorem{lemma}{Lemma}
\newtheorem{theorem}{Theorem}
\title{Non-stable subnormal contractions have nontrivial hyperinvariant subspaces}
\author{Maria F. Gamal'}
\date{Source: arXiv:2604.26044v1 (CC BY 4.0)}
\begin{document}
\maketitle

\noindent\emph{Source and licence.} Non-stable subnormal contractions have nontrivial hyperinvariant subspaces. Version arXiv:2604.26044v1 (CC BY 4.0), distributed by arXiv under the Creative Commons Attribution 4.0 International licence, http://creativecommons.org/licenses/by/4.0/, as recorded in the arXiv licence metadata for this exact version and shown on the abstract page https://arxiv.org/abs/2604.26044v1. Full licence notice: \texttt{licenses/arxiv-2604.26044-CC-BY-4.0.txt}.\par\smallskip

\noindent\textit{Editorial scope.} Counted unit: the article's theorem thm12 (source0.tex lines 679--684). The article's complete original proof of it is delivered with it (source0.tex lines 686--703, one \texttt{proof} environment of 18 lines). No mathematics is written, repaired or re-derived: the statement and the proof are the article's own lines, in the article's own notation, and no retained line is substituted or re-flowed.  Retained uncounted as the source-local context the proof needs: lines 276--284; lines 430--434.  Nothing was omitted from any retained span, so the package carries no omission record.  No retained line contains a citation, so the wrapper carries no bibliography and no bibliography entry.  No retained line contains a figure environment, an image-inclusion command or drawing code, so no image is delivered and the package declares no asset dependency.  Editorial formatting only, in addition: an AMS \texttt{amsart} preamble that restates the article's own declarations and macros used by the retained lines, namely the environments \texttt{lemma}, \texttt{theorem} on separate counters, together with the standard packages those lines need.  The retained environments print in this document as Lemma 1, Theorem 1 in order of appearance, whereas the article prints the same items with the numbering of its own full document.  The complete native source of the article as distributed in the arXiv e-print for this exact version is bundled unchanged as \texttt{sources/199357/source0.tex}.
\begin{lemma}\label{lemtau}  Suppose that $T$ is a contraction,  $\tau\subset\mathbb T$ is a Borel set,  $m(\tau)<1$, 
and $U_\tau$ is the operator of multiplication by the independent variable acting on $L^2(\tau,m)$.   
Set 
\begin{equation*}  R=U_{\tau}\oplus T.
\end{equation*}
If there exists  $\mathcal M\in\operatorname{Lat}R$ such that  $R|_{\mathcal M}\cong U_{\mathbb T}$, then 
there exists a reducing subspace $\mathcal M_0$ for $T$ such that $\mathcal M_0\neq\{0\}$ and
   $T|_{\mathcal M_0}$ is unitary. 
\end{lemma}

\begin{theorem}\label{thm11} Suppose that $T$ is a non-stable (that is, not of class $C_{0\cdot}$) subnormal contraction, 
$U$ is a unitary summand of the minimal normal extension of $T$, and $m$ is a scalar-valued spectral measure for $U$. 
Then  there exists $\mathcal M\in \operatorname{Lat}T$ such that  $T|_{\mathcal M}\cong U_{\mathbb T}$ or there exists 
a singular inner function $\theta$ such that the range of $\theta(T)$ is not dense.
\end{theorem}

\begin{theorem}\label{thm12} Suppose that $T$ is a non-stable (that is, not of class $C_{0\cdot}$) subnormal contraction, such that 
the spectral measure of a unitary summand of the minimal normal extension of $T$ is absolutely continuous 
(with respect of $m$). 
Then  there exists $\{0\}\neq\mathcal M_0\in \operatorname{Lat}T$ such that  $T|_{\mathcal M_0}$ is unitary  or there exists 
a singular inner function $\theta$ such that the range of $\theta(T)$ is not dense.
\end{theorem}

\begin{proof} By assumption, there exists  a Borel set $\tau\subset\mathbb T$  such that  $m(\tau)>0$ 
and a scalar-valued spectral measure of a unitary summand of the minimal normal extension of $T$ is $m|_{\tau}$. 
If  $m(\tau)=1$, then $T$ satisfies the assumption of Theorem \ref{thm11}. Thus, assume that $m(\tau)<1$.  
Set 
\begin{equation*}R=U_{\mathbb T\setminus\tau}\oplus T.
\end{equation*}
Then $R$ satisfies the assumption of Theorem \ref{thm11}. 
Consequently,  there exists $\mathcal M\in \operatorname{Lat}R$ such that  $R|_{\mathcal M}\cong U_{\mathbb T}$ or there exists 
a singular inner function $\theta$ such that the range of $\theta(R)$ is not dense.

If there exists $\mathcal M\in \operatorname{Lat}R$ such that  $R|_{\mathcal M}\cong U_{\mathbb T}$, 
then 
there exists $\{0\}\neq\mathcal M_0\in\operatorname{Lat}T$ such that  $T|_{\mathcal M_0}$ is unitary by Lemma \ref{lemtau} 
(applied to $\mathbb T\setminus\tau$; the inequality $m(\mathbb T\setminus\tau)<1$ follows from the assumption $m(\tau)>0$). 
If there exists 
a singular inner function $\theta$ such that the range of $\theta(R)$ is not dense, then the range of $\theta(T)$ is not dense, because  $\theta(R)=\theta(U_{\mathbb T\setminus\tau})\oplus \theta(T)$ and the range of 
$\theta(U_{\mathbb T\setminus\tau})$ is dense.
\end{proof}

\end{document}
